\begin{casebox}
\textbf{知识点小案例：旧 owner 的迟到写入为什么需要 route epoch。}
特例是 shard S 从节点 A 迁移到节点 B。客户端还拿着旧路由，把写请求发给 A；同时新请求已经到达 B。如果请求和存储记录没有 route epoch，A 的迟到写入可能覆盖新 owner 的事实。加入 epoch 后，A 只能接受自己仍是 owner 的 epoch；旧 epoch 请求必须拒绝或转发。由此推出规律：分片迁移不是移动数据那么简单，而是控制“谁有资格写入某个逻辑事实”。工程上路由表、请求、manifest 和恢复日志都要带 epoch，并写清旧请求、迟到响应和切换点。
\end{casebox}
